\documentclass[conference,9pt]{IEEEtran}
\IEEEoverridecommandlockouts
\usepackage{cite}
\usepackage{amsmath,amssymb,amsfonts}
\usepackage{algorithmic}
\usepackage{graphicx}
\usepackage{textcomp}
\usepackage{xcolor}
\usepackage[italian]{babel}

\begin{document}

\title{Canny Edge Detector: studio e implementazione}

\author{\IEEEauthorblockN{Martin Trajkovski}
\IEEEauthorblockA{\textit{Corso di Laurea in Scienze Informatiche} \\
\textit{Universit\`a degli Studi di Parma}\\
Matricola: 397464}
}

\maketitle

\begin{abstract}
Studio e implementazione da zero in NumPy del Canny Edge Detector. La pipeline
copre smoothing gaussiano, gradiente con filtri di Sobel, non-maximum
suppression, double threshold e hysteresis, secondo il programma del corso
(filtri lineari, derivate parziali, convoluzione 2D). Si discutono le scelte
implementative pi\`u rilevanti e l'effetto dei parametri $\sigma$ e delle
soglie sul risultato finale.
\end{abstract}

\section{Introduzione}
L'\textit{edge detection} mira a estrarre i contorni di un'immagine come
luoghi di forte variazione di luminanza, riducendo il dato visivo a una
rappresentazione strutturale compatta. Il rilevatore di
Canny~\cite{canny1986} combina filtraggio, derivazione e sogliatura in una
pipeline che soddisfa tre criteri di ottimalit\`a (buona detection, buona
localizzazione, risposta singola per bordo) ed \`e tuttora standard di
riferimento~\cite{gonzalez}.

\section{Pipeline}
Sia $I:\mathbb{R}^2\to\mathbb{R}$ l'immagine in scala di grigi. La pipeline
applicata \`e:

\textbf{1. Smoothing.} Convoluzione con kernel gaussiano
$G_\sigma(x,y)=\frac{1}{2\pi\sigma^2}\exp\!\left(-\frac{x^2+y^2}{2\sigma^2}\right)$
per ridurre il rumore: $I_s = I * G_\sigma$.

\textbf{2. Gradiente.} Approssimazione discreta tramite filtri di Sobel:
$G_x = I_s * S_x,\ G_y = I_s * S_y$, da cui magnitudo
$M=\sqrt{G_x^2+G_y^2}$ e direzione $\theta=\arctan2(G_y,G_x)$.

\textbf{3. Non-Maximum Suppression.} Per ogni pixel si discretizza $\theta$ in
uno dei quattro settori $\{0^\circ,45^\circ,90^\circ,135^\circ\}$ e si conserva
$M(i,j)$ solo se massimo locale lungo la direzione del gradiente.

\textbf{4. Double Threshold.} Definite due soglie $T_l<T_h$ relative a
$\max M$, i pixel sono classificati come \emph{forti} ($\ge T_h$),
\emph{deboli} ($T_l\le \cdot < T_h$) o scartati.

\textbf{5. Hysteresis.} I pixel deboli vengono promossi a edge se connessi
(8-vicinato) a un pixel forte; gli altri sono soppressi.

\section{Scelte implementative}
L'implementazione \`e in \texttt{NumPy} puro, senza ricorrere a OpenCV o
\texttt{scipy.ndimage}, per esplicitare ogni passaggio della pipeline.

\textbf{Conversione in luminanza.} Si usa la combinazione ITU-R BT.601
$Y=0.299R+0.587G+0.114B$ anzich\'e una media uniforme, perch\'e pesa i canali
in funzione della sensibilit\`a percettiva e produce un'immagine di intensit\`a
pi\`u coerente con la luminosit\`a vista.

\textbf{Convoluzione.} Implementata con
\texttt{numpy.lib.stride\_tricks.sliding\_window\_view} ed
\texttt{einsum}, che vettorializza il prodotto kernel-finestra evitando
cicli Python espliciti. Il padding \`e \texttt{edge} (replica del bordo) per
non introdurre artefatti scuri ai margini, e il kernel viene flippato in modo
da realizzare la vera convoluzione e non la semplice correlazione.

\textbf{Sobel vs Prewitt.} Si \`e scelto il filtro di Sobel perch\'e i pesi
$(1,2,1)$ sull'asse ortogonale equivalgono a uno smoothing gaussiano discreto
$3\times 3$, rendendo il gradiente meno sensibile al rumore residuo rispetto a
Prewitt (pesi uniformi $(1,1,1)$), a parit\`a di costo computazionale.

\textbf{NMS a 4 settori.} La direzione $\theta\in[0,\pi)$ viene quantizzata in
quattro settori da $45^\circ$. \`E un compromesso standard: evita
l'interpolazione bilineare tra pixel (pi\`u accurata ma pi\`u costosa) e
mantiene il confronto su vicini esatti.

\textbf{Soglie relative e hysteresis.} $T_l$ e $T_h$ sono frazioni di
$\max M$, cos\`i gli stessi parametri funzionano su range di intensit\`a
diversi; l'hysteresis itera fino a punto fisso, esplicitando la propagazione
di connettivit\`a in 8-vicinato.

\section{Risultati}
L'algoritmo \`e stato testato su immagini reali con $\sigma=1.4$,
kernel $5\times 5$, $T_l=0.05\cdot\max M$, $T_h=0.15\cdot\max M$.
La Figura~\ref{fig:pipeline} mostra gli stadi intermedi: lo smoothing
attenua la tessitura ma preserva i contorni principali, la magnitudo del
gradiente li evidenzia con spessore di alcuni pixel, l'NMS li assottiglia
a un pixel e la sogliatura con hysteresis rimuove le risposte spurie
mantenendo le strutture connesse. Esperimenti sul parametro $\sigma$
($0.8$, $1.4$, $2.5$) confermano il \textit{trade-off} tra dettaglio e
soppressione del rumore tipico del rilevatore.

\begin{figure}[h]
\centering
\includegraphics[width=\columnwidth]{pipeline.png}
\caption{Stadi della pipeline: grayscale, smoothing, magnitudo del gradiente,
NMS, double threshold, edge finali.}
\label{fig:pipeline}
\end{figure}

\section{Conclusioni}
Le quattro idee chiave del corso, ovvero convoluzione 2D, derivate via Sobel,
sogliatura adattiva e connessione di componenti, bastano a costruire un
rilevatore di bordi competitivo, con $\sigma$ come parametro pi\`u critico.

\begin{thebibliography}{1}
\bibitem{canny1986} J. Canny, ``A Computational Approach to Edge Detection,''
\textit{IEEE Trans. Pattern Anal. Mach. Intell.}, vol. PAMI-8, no. 6, 1986.
\bibitem{gonzalez} R. C. Gonzalez and R. E. Woods, \textit{Digital Image
Processing}, 4th ed., Pearson, 2018.
\end{thebibliography}

\end{document}
